\documentclass[12pt,a4paper]{article}
\usepackage[utf8]{inputenc}
\usepackage[T1]{fontenc}
\usepackage{lmodern}
\usepackage[french]{babel}
\usepackage[a4paper,margin=2.3cm]{geometry}
\usepackage{amsmath,amssymb}
\usepackage{enumitem}
\setlist[itemize]{label=\textbullet}
\usepackage[most]{tcolorbox}
\usepackage{fancyhdr}

\pagestyle{fancy}
\fancyhf{}
\fancyhead[L]{\small Exercice de maison}
\fancyhead[R]{\small Suites arithmétiques}
\fancyfoot[C]{\small\thepage}
\renewcommand{\headrulewidth}{0.4pt}

\setlength{\parindent}{0pt}
\setlength{\parskip}{6pt}

\newtcolorbox{resultat}{colback=gray!8,colframe=black,boxrule=0.5pt,arc=1pt,
  left=6pt,right=6pt,top=3pt,bottom=3pt}

\begin{document}

\begin{center}
{\LARGE\bfseries Exercice de maison}\\[4pt]
{\large Suite arithmétique}
\end{center}

\vspace{4pt}
\noindent Nom et prénom : \dotfill \quad Classe : \dotfill

\vspace{6pt}
\begin{tcolorbox}[colback=white,colframe=black,boxrule=0.5pt,arc=1pt,title=\textbf{Énoncé},
  coltitle=black,colbacktitle=gray!15]
On donne la suite $(U_n)$ définie par :
\[
\begin{cases}
U_0 = 3\\
U_{n+1} = U_n - 7
\end{cases}
\]
\begin{enumerate}[nosep]
  \item Calculer $U_1$ et $U_2$.
  \item Montrer que $(U_n)$ est une suite arithmétique dont on précisera la raison et le premier terme.
  \item Exprimer $U_n$ en fonction de $n$.
  \item Donner le sens de variation de $(U_n)$.
  \item Calculer $S = U_0 + U_1 + \dots + U_{100}$.
\end{enumerate}
\end{tcolorbox}

\section*{Solution}

\subsection*{1. Calcul de $U_1$ et $U_2$}
On utilise la relation de récurrence $U_{n+1} = U_n - 7$.

Pour calculer $U_1$, on remplace $n$ par $0$, on obtient :
\quad $U_{0+1} = U_1 = U_0 - 7 = 3 - 7 = -4$.

Pour calculer $U_2$, on remplace $n$ par $1$, on obtient :
\quad $U_{1+1} = U_2 = U_1 - 7 = -4 - 7 = -11$.

\begin{resultat}
$U_1 = -4$ \qquad et \qquad $U_2 = -11$.
\end{resultat}

\subsection*{2. Nature de la suite $(U_n)$}
\textbf{Rappel du cours (définition).} Une suite numérique $(U_n)_{n\in\mathbb{N}}$ est dite arithmétique
si, $U_0$ étant donné, il existe un nombre réel $r$ tel que pour tout entier naturel $n$ :
$U_{n+1} = U_n + r$. Le réel $r$ est appelé la raison. Autrement dit, $(U_n)$ est arithmétique de
raison $r$ lorsque pour tout entier naturel $n$ : $U_{n+1} - U_n = r$.

Pour tout entier naturel $n$, on a :
\[
U_{n+1} = U_n - 7 = U_n + (-7),
\qquad \text{c'est-à-dire} \qquad
U_{n+1} - U_n = (U_n - 7) - U_n = -7.
\]
La différence entre deux termes consécutifs est égale à la constante $-7$, qui ne dépend pas de $n$.
D'après la définition, la suite $(U_n)$ est donc une suite arithmétique.

\begin{resultat}
$(U_n)$ est la suite arithmétique de raison $r = -7$ et de premier terme $U_0 = 3$.
\end{resultat}

\subsection*{3. Expression de $U_n$ en fonction de $n$}
\textbf{Rappel du cours (propriété).} Si $(U_n)_{n\in\mathbb{N}}$ est la suite arithmétique de premier
terme $U_0$ et de raison $r$, alors pour tout entier naturel $n$ : $U_n = U_0 + n\,r$.
Plus généralement, pour tous entiers naturels $n$ et $p$ : $U_n = U_p + (n - p)\,r$.

Ici $U_0 = 3$ et $r = -7$. En appliquant la propriété avec $p = 0$, on a pour tout entier naturel $n$ :
\[
U_n = U_0 + (n - 0)\,r = 3 + n \times (-7) = 3 - 7n.
\]
\emph{Vérification :} $U_1 = 3 - 7 \times 1 = -4$ et $U_2 = 3 - 7 \times 2 = -11$, ce qui correspond
aux résultats de la question 1.

\begin{resultat}
Pour tout entier naturel $n$ : \quad $U_n = 3 - 7n$.
\end{resultat}

\subsection*{4. Sens de variation de $(U_n)$}
\textbf{Rappel du cours (propriété).}
\begin{itemize}[nosep]
  \item Une suite arithmétique est croissante si sa raison est positive.
  \item Une suite arithmétique est décroissante si sa raison est négative.
  \item Une suite arithmétique est constante si sa raison est nulle.
\end{itemize}

La suite $(U_n)$ est arithmétique de raison $r = -7$. Comme $-7 < 0$, sa raison est négative.
On le retrouve directement : pour tout entier naturel $n$, $U_{n+1} - U_n = -7 < 0$, donc $U_{n+1} < U_n$.

\begin{resultat}
Comme $r = -7$ et $-7 < 0$, alors la suite $(U_n)$ est décroissante (strictement, car $r \neq 0$).
\end{resultat}

\subsection*{5. Calcul de $S = U_0 + U_1 + \dots + U_{100}$}
\textbf{Rappel du cours (propriété).} Si $(U_n)_{n\in\mathbb{N}}$ est une suite arithmétique, alors pour
tout entier naturel $n$ :
\[
U_0 + U_1 + \dots + U_n = (n + 1) \times \frac{U_0 + U_n}{2}.
\]
Autrement dit, la somme de termes consécutifs d'une suite arithmétique est égale à :
\[
S = (\text{nombre de termes}) \times \frac{\text{premier terme} + \text{dernier terme}}{2}.
\]

\begin{itemize}[nosep]
  \item Nombre de termes : de $U_0$ à $U_{100}$, il y a $100 - 0 + 1 = 101$ termes.
  \item Premier terme : $U_0 = 3$.
  \item Dernier terme : d'après la question 3, $U_{100} = 3 - 7 \times 100 = 3 - 700 = -697$.
\end{itemize}

On applique la propriété avec $n = 100$ :
\[
S = (100 + 1) \times \frac{U_0 + U_{100}}{2}
  = 101 \times \frac{3 + (-697)}{2}
  = 101 \times \frac{-694}{2}
  = 101 \times (-347)
  = -35\,047.
\]

\begin{resultat}
$S = U_0 + U_1 + \dots + U_{100} = -35\,047$.
\end{resultat}

\end{document}
